\section{Introduction}
\label{sec:intro}

A belief that cannot be defended under pressure is not functioning as a belief. An
assistant that folds to a fabricated statistic is unusable for research support or
review; an ensemble of agents that converges under social pressure rather than evidence
collapses to consensus instead of staying adversarial; an agent that abandons a \emph{correct}
position under authority pressure fails in the direction hardest to detect, because the
surface behavior --- politeness, agreement, a plausible-sounding concession --- looks
like good conversational hygiene. Persistent conditioning of a language model's stance
or persona is nonetheless almost universally implemented by placing it in the context
window, alongside the interlocutor's own content, arbitrated by the model's own chain of
thought. The conditioning and the pressure that will test it share a channel.

\textbf{We show this dissociates.} Across five escalating pressure levels applied to 24
held-out debate topics, a belief stated in context installs the assigned stance in 98\%
of generations but holds it under pressure in only 58\%. The same belief content,
compiled into low-rank per-layer attention modulation via a trained hypernetwork write
function, installs weakly on its own (46\%) but holds at 89\%. Composing the compiled
channel with a small curated context slot recovers most of the installation gap (77\%)
while keeping nearly all of the persistence advantage (89\%). Installation and
persistence are not two measurements of one capability; they are separable properties,
and the channel that installs best is not the channel that holds best.

The gap is not a uniform lift. It concentrates precisely where context conditioning
collapses hardest --- logical counterargument (mean conviction 2.93 for context-only,
4.78 for compiled) and emotional pressure (1.99 vs.\ 3.82) --- which is a mechanism
signature, not an artifact of an easy benchmark.

\textbf{Why.} Text placed in context is read, reasoned about, and is available to be
argued against; instruction tuning specifically rewards accommodating an articulate
objection. A compiled modulation is not read as an assertion in that sense --- it
conditions the forward pass rather than occupying a position in it, and it is not
present in the chain-of-thought the model reasons over (Section~\ref{sec:why-channels-differ}).

\textbf{Alternative explanation 1: is this just activation steering?} Contrastive
activation vectors extracted from identical belief content, tuned honestly on held-out
training topics, install at 42\% --- matching compiled modulation's 46\% --- and then
collapse to 14\% hold under the same pressure protocol where compiled modulation holds
89\%. This is not an artifact of a weak baseline: an independent systematic stress test
of activation steering under adversarial perturbation reports directional robustness
dropping by up to 64 percentage points and post-attack confidence collapsing to or below
0.25 across four extraction methods and five models from 1.5B to 30B parameters
\citep{le2026steeringrobustness}. Section~\ref{sec:isolation} reports the full ablation.

\textbf{Alternative explanation 2: is this just prompt compression?} Compressed tokens
still occupy context positions and still dilute against a growing transcript; compiled
tensors occupy none. We return to this distinction, and to what remains open about it,
in Section~\ref{sec:isolation}.

\textbf{Boundary.} The mechanism carries disposition, not facts. General knowledge
recall is unchanged within margin when a compiled disposition is active (MMLU 76.5\% vs.\
74.8\%, paired McNemar \(p=0.11\)), and specific, unfamiliar content --- fabricated
statistics, unfamiliar proper nouns, exact action strings --- does not survive the
compilation bottleneck and is deliberately routed to a curated context slot instead
(Section~\ref{sec:boundary}).

\textbf{Discrimination, not rigidity.} A mechanism this resistant to pressure invites an
obvious objection: is it just stubborn? We measure this directly. Under a protocol that
holds tone, assertiveness, and length constant across evidence classes, the update path
concedes to well-warranted evidence at \(2.4\times\) the rate it concedes to facially
deficient evidence (42.5\% vs.\ 17.5\%, \(\Delta = +25.0\) points), and concedes 0\% of
the time to the four contentless pressure types (Section~\ref{sec:discrimination}). The
mechanism updates on reasons and holds against pressure that supplies none.

\paragraph{Scope.} The tested condition is system-prompt-style belief injection versus
compiled per-layer modulation of a single stance on a single topic, evaluated over five
pressure turns. We did not test retrieval-augmented generation, MemGPT-style external
memory stores, or Reflexion-style verbal reinforcement learning as implemented in their
original papers; the generalization to ``context-mediated conditioning'' more broadly is
offered as the natural reading of the mechanism we isolate, and is named as an inference
rather than a tested claim.

\paragraph{Contributions.} (1) We report a dissociation between installation and
persistence across two conditioning channels for the same belief content, isolated by
construction (Section~\ref{sec:dissociation}). (2) We isolate the effect against the
two most obvious competing explanations --- fixed-direction activation steering and
prompt compression --- and against a frozen-offset ablation of the mechanism itself,
localizing what does the work (Section~\ref{sec:isolation}). (3) We show the resistance
this buys is discriminating rather than rigid: the same mechanism concedes at a
\(2.4\times\) higher rate to well-warranted evidence than to facially deficient pressure
(Section~\ref{sec:discrimination}). (4) We connect the result to debate-based scalable
oversight, whose soundness guarantee assumes debaters hold positions for reasons ---
an assumption every existing remedy places in the same channel the opponent argues in
(Section~\ref{sec:related}).
